\section*{الفصل \ref{ch:g8:ice-on-earth} --- الجليد على الأرض}

\begin{solution}{exo:g8:ice-on-earth:1}
(أ) منطقة التراكم. (ب) منطقة الذوبان. (ج) المقدمة. (د) \omterm{def:g8:ice-on-earth:glacial-landforms}{ركام
جليدي} (طرفي).
\end{solution}

\begin{solution}{exo:g8:ice-on-earth:2}
(أ) خطأ: يتكون \omterm{def:g8:ice-on-earth:glacier}{النهر الجليدي} من \omterm{def:g4:clouds-rain-snow:precipitation}{ثلج} سقط على اليابسة وتحول إلى جليد؛ أما ماء
\omterm{def:g2:water-on-earth:ocean}{البحر} المتجمد فهو \omterm{def:g8:ice-on-earth:sea-ice}{الجليد البحري}. (ب) صواب. (ج) خطأ: جليده لا يزال يجري نحو
الأسفل؛ وتتراجع المقدمة لأن الجليد يذوب أسرع مما يجري. (د)
صواب.
\end{solution}

\begin{solution}{exo:g8:ice-on-earth:3}
(أ) \omterm{def:g8:ice-on-earth:glacial-landforms}{وادٍ على شكل حرف U}. (ب) صخرة ضالة. (ج) \omterm{def:g8:ice-on-earth:ice-sheet}{رف جليدي}. (د) \omterm{def:g8:ice-on-earth:permafrost}{التربة الصقيعية}.
\end{solution}

\begin{solution}{exo:g8:ice-on-earth:4}
(أ) القارة القطبية الجنوبية وغرينلاند. (ب) القارة القطبية الجنوبية، بنحو
\qty{58}{m} (غرينلاند: نحو \qty{7.4}{m}). (ج) نحو العُشر (10~\%).
\end{solution}

\begin{solution}{exo:g8:ice-on-earth:5}
(أ) نحو \num{7.05}~مليون~\unit{km^2} عام 1979 و\num{3.57}~مليون~\unit{km^2}
عام 2012. (ب) $7.05 - 3.57 = 3.48$: أي نحو \num{3.5}~مليون~\unit{km^2}.
(ج) $3.48 \div 7.05 \approx 0.49$: أي نحو النصف. (د) السنة الواحدة تتوقف
على \omterm{def:g1:watching-the-weather:weather}{طقس} ذلك الصيف: فعام 2012 يقع تحت الخط بكثير. أما خط الاتجاه فيستعمل
السنوات كلها، فلا يغيره صيف استثنائي واحد إلا قليلًا جدًا.
\end{solution}

\begin{solution}{exo:g8:ice-on-earth:6}
(أ) $2004 - 1941 = 63$ سنة. (ب) $12000 \div 63 \approx 190.5$: أي نحو
\qty{190}{m} في السنة. (ج) $2021 - 1984 = 37$ سنة؛ $3200 \div 37 \approx
86$: أي نحو \qty{86}{m} في السنة. (د) \omterm{def:g1:landscapes-around-us:water}{نهر} ميوير الجليدي، بأكثر من ضعف السرعة.
\end{solution}

\begin{solution}{exo:g8:ice-on-earth:7}
في \emph{منطقة التراكم}، عاليًا على \omterm{def:g1:landscapes-around-us:landform}{الجبل}، يسقط كل سنة من \omterm{def:g4:clouds-rain-snow:precipitation}{الثلج} أكثر مما
يذوب. فيدفن \omterm{def:g4:clouds-rain-snow:precipitation}{الثلج} الجديد \omterm{def:g4:clouds-rain-snow:precipitation}{الثلج} القديم ويضغطه، فيتحول بعد نحو سنة إلى
\emph{\omterm{def:g4:clouds-rain-snow:precipitation}{ثلج} حبيبي}، ثم، في عمق أكبر وبعد زمن أطول بكثير، إلى جليد صلب. وتحت
وزنه \emph{يجري} الجليد ببطء نزولًا في \omterm{def:g1:landscapes-around-us:landform}{الوادي}. وفي الأسفل، في \emph{منطقة
الذوبان}، يذوب أسرع مما يسقط \omterm{def:g4:clouds-rain-snow:precipitation}{الثلج}؛ والجليد الذي يبلغ المقدمة جليد تكوّن في
الأعلى، وكثيرًا ما تكوّن منذ زمن بعيد.
\end{solution}

\begin{solution}{exo:g8:ice-on-earth:8}
(أ) في الكأس B: فالجليد الذائب يسيل إلى الماء ويرفعه؛ وفي الكأس A يبقى
المستوى حيث كان. (ب) الكأس A هي \omterm{def:g8:ice-on-earth:ice-sheet}{الجبل الجليدي} (جليد يطفو في \omterm{def:g2:water-on-earth:ocean}{البحر} أصلًا)؛
والكأس B هي \omterm{def:g8:ice-on-earth:ice-sheet}{الغطاء الجليدي} (جليد على اليابسة). (ج) \omterm{def:g8:ice-on-earth:ice-sheet}{الرف الجليدي} يحجز \omterm{def:g8:ice-on-earth:glacier}{الأنهار
الجليدية} التي خلفه. فحين يتفكك، تتسارع هذه \omterm{def:g8:ice-on-earth:glacier}{الأنهار الجليدية} وتصب في \omterm{def:g2:water-on-earth:ocean}{البحر}
مزيدًا من جليد اليابسة، وهذا يرفعه.
\end{solution}

\begin{solution}{exo:g8:ice-on-earth:9}
(أ) \omterm{def:g8:ice-on-earth:glacial-landforms}{وادٍ على شكل حرف U}، نحته \omterm{def:g8:ice-on-earth:glacier}{نهر جليدي}. (ب) حمله \omterm{def:g8:ice-on-earth:glacier}{نهر جليدي} من موضع بعيد في
أعلى \omterm{def:g1:landscapes-around-us:landform}{الوادي}، حيث يوجد الغرانيت، وتركه حين ذاب الجليد: إنه صخرة ضالة. (ج) ركام
جليدي جانبي: يبين ارتفاعه كم كان سُمك \omterm{def:g8:ice-on-earth:glacier}{النهر الجليدي} يومًا؛ فقد بلغ الجليد هذا
الارتفاع على الأقل على جانب \omterm{def:g1:landscapes-around-us:landform}{الوادي}.
\end{solution}

\begin{solution}{exo:g8:ice-on-earth:10}
(أ) $7.4 \div 10 = \qty{0.74}{m}$. (ب) $58 \div 100 = \qty{0.58}{m}$.
(ج) $58 \div 7.4 \approx 7.8$: أي نحو 8 مرات.
\end{solution}

\begin{solution}{exo:g8:ice-on-earth:11}
(أ) $7.21 \times 10^6 = \qty{7210000}{km^2}$ و$4.65 \times 10^6 =
\qty{4650000}{km^2}$. (ب) $7.21 - 4.65 = \num{2.56}\text{ مليون }\unit{km^2}$، أي
\qty{2560000}{km^2}. (ج) $2.56 \div 3 \approx 0.85$: أي نحو
\num{0.85}~مليون~\unit{km^2} في كل عقد، وهو أكثر بقليل من انخفاض خط الاتجاه
البالغ \num{0.76}~مليون~\unit{km^2}؛ فمتوسط كل عقد لا يزال يحمل صعود وهبوط
أصيافه العشرة، لكن طريقتي قراءة المخطط كلتيهما تبينان الخسارة الكبيرة نفسها.
(د) لا يمكن أن ينخفض الامتداد إلى ما دون الصفر، وخط الاتجاه لا يصف إلا
السنوات التي قيست: وما يحدث بعد ذلك يتوقف على كيفية تغير \omterm{def:g5:weather-and-climate:climate}{مناخ} المنطقة القطبية
الشمالية، فالخط ليس تنبؤًا.
\end{solution}

\begin{solution}{exo:g8:ice-on-earth:12}
(أ) $120 \div 1 = 120$: كان \omterm{def:g1:landscapes-around-us:water}{الساحل} أبعد في عرض \omterm{def:g2:water-on-earth:ocean}{البحر} بنحو \qty{120}{km}.
(ب) كان قد تبخر من \omterm{def:g2:water-on-earth:ocean}{المحيط}، وسقط ثلجًا على \omterm{def:g3:continents-and-oceans:continent}{القارات} وبقي هناك جليدًا في \omterm{def:g8:ice-on-earth:ice-sheet}{الأغطية
الجليدية} الكبرى. (ج) يتكون \omterm{def:g8:ice-on-earth:glacial-landforms}{الركام الجليدي} الطرفي عند نهاية \omterm{def:g8:ice-on-earth:glacier}{نهر جليدي}. فوجوده
في \omterm{def:g1:landscapes-around-us:landform}{سهل} بعيد عن أي \omterm{def:g1:landscapes-around-us:landform}{جبال} يبين أن غطاءً جليديًا بلغ ذلك \omterm{def:g1:landscapes-around-us:landform}{السهل} يومًا، في زمن كان
الجليد يغطي فيه من اليابسة أكثر بكثير مما يغطيه اليوم.
\end{solution}

\begin{solution}{pb:g8:ice-on-earth:1}
\textbf{1.} $1080 - 0 = \qty{1080}{m}$.

\textbf{2.} \babelsublr{1960--1970}: \qty{180}{m}؛ \babelsublr{1970--1980}: \qty{150}{m}؛
1980--1990: \qty{190}{m}؛ 1990--2000: \qty{180}{m}؛ 2000--2010:
\qty{180}{m}؛ 2010--2020: \qty{200}{m}. أقل تراجع في 1970--1980، وأكبره
في 2010--2020.

\textbf{3.} لا: \omterm{def:g8:ice-on-earth:glacier}{فالنهر الجليدي} المتراجع لا يزال يجري نحو الأسفل. وتتراجع
مقدمته لأن الجليد يذوب أسرع مما يصل.

\textbf{4.} ركام جليدي طرفي: حمل الجليد الحطام الصخري نزولًا في \omterm{def:g1:landscapes-around-us:landform}{الوادي}
وألقاه عند مقدمته، حيث كان الجليد يذوب، طوال بقاء المقدمة
هناك.

\textbf{5.} $1080 \div 6 = \qty{180}{m}$ في كل عقد.

\textbf{6.} $180 \div 10 = \qty{18}{m}$ في السنة.

\textbf{7.} $190 \div 18 \approx 10.6$: أي أسرع بنحو 11 مرة.

\textbf{8.} أقل: فقد خسر في منطقة ذوبانه من الجليد أكثر مما كسب في منطقة
تراكمه، فتراجعت مقدمته.

\textbf{9.} على شكل حرف U: قاع عريض وجوانب شديدة الانحدار، لأن \omterm{def:g8:ice-on-earth:glacier}{النهر
الجليدي} وسّع \omterm{def:g1:landscapes-around-us:landform}{الوادي} وعمّقه حين كان يملؤه.

\textbf{10.} $1980 - 1080 = \qty{900}{m}$.

\textbf{11.} $900 \div 18 = 50$ سنة: أي في عام \textbf{2070}.

\textbf{12.} تغيرت السرعة من عقد إلى عقد (من 150 إلى \qty{200}{m})،
والتراجع المقبل يتوقف على تساقط \omterm{def:g4:clouds-rain-snow:precipitation}{الثلج} ودرجات حرارة الصيف في السنوات القادمة:
فالمعدل الماضي ليس تنبؤًا.
\end{solution}
